% Generated draft: requires mathematical and source review.
\documentclass[11pt]{article}
\usepackage[margin=25mm]{geometry}
\usepackage{fontspec}
\setmainfont{Latin Modern Roman}
\usepackage{amsmath,amssymb,mathtools,unicode-math}
\setmathfont{Latin Modern Math}
\usepackage{enumitem,needspace,xurl,hyperref}
\setlength{\parindent}{0pt}
\setlength{\parskip}{5pt}
\setlength{\emergencystretch}{3em}
\newcommand{\N}{\mathbb{N}}
\newcommand{\Z}{\mathbb{Z}}
\newcommand{\Q}{\mathbb{Q}}
\newcommand{\R}{\mathbb{R}}
\newcommand{\C}{\mathbb{C}}
\newcommand{\F}{\mathbb{F}}
\newcommand{\A}{\mathbb{A}}
\newcommand{\PP}{\mathbb{P}}
\newcommand{\E}{\mathbb{E}}
\newcommand{\eps}{\varepsilon}
\newcommand{\dd}{\,\mathrm{d}}
\begin{document}
\section*{UnsolvedMath ID: 10400102}
\textbf{Source problem code:} catalogue entry
\subsection{Definitions and mathematical statement}
A quandle \((Q,*)\) satisfies \(x*x=x\), each right translation \(x\mapsto x*y\) is bijective, and \((x*y)*z=(x*z)*(y*z)\). It is connected if the group generated by right translations acts transitively. For an abelian coefficient group \(A\), written multiplicatively, a normalized quandle 2-cocycle \(\theta\) satisfies
\[
 \theta(x,x)=1,\qquad
 \theta(x,y)\theta(x*y,z)=\theta(x,z)\theta(x*z,y*z).
\]
Let \(\alpha\in H_Q^2(Q;A)\) be its cohomology class.

In the standard finite-quandle convention, color an oriented knot diagram by \(Q\) using the quandle crossing rule, and define
\[
 \Phi_\alpha(K)=\sum_{\mathcal C\text{ a }Q\text{-coloring}}
 \prod_{c\text{ a crossing}}\theta(a_c,b_c)^{\varepsilon_c}
 \quad\in\mathbb Z[A].
\]
Here \(b_c\) is the over-arc color, \(a_c\) the source under-arc color in the chosen coorientation convention, and \(\varepsilon_c\) the crossing sign. This requires a consistent convention at negative crossings; the excerpt does not reproduce it.

The source asks to compute these cocycle invariants for each knot and the intended second cohomology classes of connected quandles. It does not specify a particular quandle, coefficient group, knot family or output form. Finiteness is not stated: an infinite connected quandle may yield infinitely many colorings, so the displayed ordinary group-ring state sum needs an additional restriction or definition in that setting. For finite input data the defining sum is enumerable; the source's computation program must therefore be distinguished from a new claim of basic computability. No closed formula, efficiency target or complete table scope is supplied.

The summary reports many special-knot/quandle computations but gives no examples or citations. Its partial-progress label remains source information pending current primary-literature review.
\subsection{Short English statement}
Compute the quandle 2-cocycle state-sum invariant for the intended knots, connected quandles and cohomology classes. The coefficient group, finite/infinite convention and desired computational output are missing, so no universal formula or new computability claim is inferred.
\subsection{Sources}
\begin{itemize}
\item Local registry: \url{data/unsolvedmath/problems.json}, permanent UnsolvedMath ID \(10400102\), source code catalogue entry; supplied statement, background, and research summary, updated 2026-08-06.
\item Ohtsuki, \emph{Problems on invariants of knots and 3-manifolds} (2002), Problem 5.6, PDF page 89: \url{https://msp.org/gtm/2002/04/gtm-2002-04-024s.pdf}. Supplied PDF-font-structure extraction accessed 2026-07-29; provenance distinguishes 231 numbered items.
\item No separate special-case computation or current resolution reference is supplied.
\end{itemize}
\end{document}
